% ======================================================================
\section{Problem Statement and Information Sets}
\label{sec:problem}
Let $\vv{x}_t\in\R^{C}$ be the vector observed at hour $t$, whose channels are the system load
$x_{t,c^\star}$ together with weather and renewable-generation series. Each step carries a UTC
timestamp $\tau_t$ and a local civil timestamp $\tau^{\ell}_t$, which differ by the zone offset
and the daylight-saving shift. Given a lookback $\mm{X}=[\vv{x}_{t-L+1},\dots,\vv{x}_t]$ and the
horizon timestamps, the task is to forecast the load $y_{t+1},\dots,y_{t+H}$ for
$H\in\{24,96,336\}$ hours, and, in the operator protocol of Section~\ref{sec:dayahead}, for
the 24 hours of day $D$ from an origin at 11:00 UTC on day $D-1$.

\subsection{Information regimes}
\label{sec:regimes}
What a forecaster is allowed to know is fixed by the regime, not by its architecture:
\begin{itemize}
\item \textbf{R0} --- the past of every channel and nothing else.
\item \textbf{R1} --- R0 plus a \emph{calendar block} $\vv{k}_t\in\R^{8}$ over lookback and
      horizon: sine and cosine of local hour of day and day of week, a weekend indicator, and
      three national-holiday indicators (public holiday, bridge day between a holiday and a
      weekend, and the 24~December--6~January period).
\item \textbf{R2} --- R1 plus a \emph{weather block}: the exogenous channels over the lookback
      as observed, and over the horizon as a day-ahead forecast would supply them, obtained by
      degrading the reanalysis with a lead-time-dependent error matched to operational skill:
      for weather variables an absolute error (\nv{proto.nwp.temp_rmse}\,K for 2\,m temperature
      at 24\,h, growing as the square root of lead time), and for wind and solar generation,
      whose forecast error scales with the installed fleet, an error relative to each system's
      capacity (\nv{proto.nwp.wind_cap} and \nv{proto.nwp.solar_cap} of the training-split
      \nv{proto.nwp.cap_pctl}th percentile at 24\,h), the normalised errors typical of national day-ahead
      renewable forecasts.
\end{itemize}
Every model in a regime receives the identical block $\mm{K}\in\R^{(L+H)\times E}$
($E=0$, $8$, or $8+C-1$), identical up to the random draw of the forecast error, which is
generated afresh for every run. A model with a native known-future interface consumes it
directly; a model without one receives it through the shared exogenous head of
Section~\ref{sec:head}; the two classical references consume the subset their specification
admits (Section~\ref{sec:baselines}). R2 is the headline regime, under which every model is
run; R0 and R1, run for the headline models and the trees, quantify what each information
source is worth to each architecture.

\subsection{Scoring}
\label{sec:metrics}
Accuracy is scored on the load channel in physical units over every test window at unit
stride and every horizon step: WAPE $=\sum|e|/\sum|y|$ over those cells, in per cent (preferred
to MAPE, which is unstable wherever the target approaches zero \cite{hyndman2006accuracy}; it
weights every hour by its load and is the metric of the tables and the text), MAE and RMSE in
megawatts, \emph{peak MAPE} (error in the height of each day's maximum, in per cent of the true
maximum), \emph{peak timing error} (hours between forecast and actual peak) and \emph{ramp MAE}
(error in successive differences). Mean squared error on the standardised target (training mean
and standard deviation removed) is kept in every record and reported for the multivariate
benchmark of Section~\ref{sec:bench}, where it is the convention; it ranks the models as WAPE
does in almost every cell. Probabilistic forecasts are scored by the mean pinball loss over the
emitted levels, the quantile-weighted CRPS approximation (the trapezoidal integral of twice the
pinball loss over the levels, which does not score the tails beyond the outer levels and is not
the CRPS of a continuous distribution), and, for the central 90\,\% interval, coverage, mean
width and the Winkler score \cite{winkler1972,gneiting2007scoring}. Every headline cell is the
mean over three training seeds; its standard deviation is reported beside it.

\subsection{Statistical comparison}
\label{sec:stats}
Pairwise differences use the Diebold--Mariano test on per-window absolute-error differentials
(seed means) with a Newey--West variance and the Harvey--Leybourne--Newbold correction
\cite{diebold1995dm,harvey1997dm}; because consecutive test windows overlap, the Bartlett
truncation lag is the larger of $n^{1/3}$ and $h-1$, the standard choice for $h$-step-ahead
errors, capped at a quarter of the sample. The overlap leaves the differential strongly
autocorrelated, and the effective sample size it rests on, $n/(1+2\sum_k w_k\rho_k)$ with the
same kernel, is far below the window count: its median over the cells is \nv{dm.n_eff_median}
and its minimum \nv{dm.n_eff_min}. Every test is therefore also run with a
stationary-bootstrap standard error of the mean differential \cite{politis1994stationary}, with
the expected block length the larger of $n^{1/3}$ and $h$ (\nv{dm.B} resamples; the block scheme of the confidence set below, which itself uses
\nv{mcs.B}) and a normal reference, and Section~\ref{sec:grids} quotes both counts beside each
other; a verdict is further counted as confirmed per seed when every seed's own test has the
same sign at $p<0.05$, and that count is quoted there too. Within each system the
$p$-values are Holm-adjusted across methods and horizons. For every (system, horizon) cell we
also report the Model Confidence Set \cite{hansen2011mcs} at level $0.10$, computed with the
same stationary bootstrap on the seed-mean per-window losses with the $T_{\max}$ statistic:
the set of models that cannot be distinguished from the best. Membership, not rank, is the
claim we make, and a set that retains many models says that the cell cannot separate them at
the available power, which at two weeks is low; Section~\ref{sec:grids} reports the set sizes by horizon for that reason.
